% =====================================================================
%  Section 2 -- Related Work
% =====================================================================

\section{Related Work}
\label{sec:related}

\subsection{Deep Learning for Flood Segmentation}
\label{sec:rw-flood}

Convolutional encoder--decoder architectures dominate the deep-learning flood-segmentation literature. Pech-May \emph{et al.}~\cite{pechmay2023flooded} applied U-Net to Sentinel-1 SAR imagery for flooded-area detection. Bahrami and Arbabkhah~\cite{bahrami2024enhanced} compared SegNet, UNet, and FCN32 for floodplain segmentation, reporting SegNet's max-pooling indices as advantageous on smooth water-surface transitions. Patil \emph{et al.}~\cite{patil2024flood} used an ensemble of ResNet34, InceptionV3, and VGG16 on Sentinel-1 SAR for flood-damage assessment, reaching an IoU of 0.758. Ghosh \emph{et al.}~\cite{ghosh2024automatic} used a Nested U-Net with FPN and EfficientNet-B7 backbone on the NASA IEEE GRSS benchmark, reporting that models trained on specific events generalise to new geographies. Chamatidis \emph{et al.}~\cite{chamatidis2024vision} introduced Vision Transformers (ViT) to the task and reported a 9--15\% improvement over comparable CNN baselines on Sentinel-1 and Sentinel-2 imagery. Most recently, Karc{\i} \emph{et al.}~\cite{karci2026deep} performed the most comprehensive systematic comparison to date, evaluating UNet, SegNet, and DeepLabV3$+$ each with EfficientNet and ResNet backbones, and concluded that UNet $+$ EfficientNet is the strongest configuration on FSSD.

Across this entire body of work, model efficiency is essentially absent from the evaluation. Hern\'andez \emph{et al.}~\cite{hernandez2022flood} gesture toward edge deployment of UAV-based flood detection, but retain heavy backbones and report no parameter, FLOP, or per-platform latency numbers. To the best of our knowledge, no prior flood-segmentation paper has measured on-device inference latency on Apple Silicon or in a browser WebAssembly sandbox, nor has any prior work quantified the per-architecture stability of INT8 post-training quantization for this task.

\subsection{Knowledge Distillation for Semantic Segmentation}
\label{sec:rw-kd}

Knowledge distillation (KD), introduced by Hinton \emph{et al.}~\cite{hinton2015distilling}, trains a small student network to match the soft outputs of a larger teacher. Romero \emph{et al.}~\cite{romero2015fitnets} extended this to feature-space alignment with FitNets, training the student to match intermediate feature representations through a learned adapter. For dense-prediction tasks, Liu \emph{et al.}~\cite{liu2019structured} proposed structured-knowledge distillation (SKD) that aligns pairwise pixel similarities; Wang \emph{et al.}~\cite{wang2020intra} proposed inter- and intra-class distance distillation specifically for segmentation; and He \emph{et al.}~\cite{he2021channel} introduced channel-wise KD for compact segmentation models. These methods routinely achieve compression ratios of $5$--$20\times$ with sub-1\% mIoU degradation on Cityscapes and ADE20K. To our knowledge, no prior work has applied KD to flood segmentation.

Two findings from the broader KD literature are particularly relevant here. First, KD gains shrink sharply as student capacity and pretraining quality grow~\cite{romero2015fitnets,beyer2022patient}. Second, on saturated benchmarks where the student matches or exceeds the teacher without distillation, the response-based term reduces to a noise-injection regulariser rather than a genuine transfer signal. Our empirical results on FSSD (Section~\ref{sec:results-kd}) are consistent with both observations.

\subsection{Lightweight Backbones for Mobile Vision}
\label{sec:rw-backbones}

MobileNetV3~\cite{howard2019searching} uses neural architecture search and the H-swish activation to produce CPU-friendly classification networks; the Small variant has approximately 2.5\,M parameters. EfficientNet-Lite~\cite{tan2019efficientnet} is a quantization-friendly redesign of EfficientNet that removes squeeze-and-excitation and Swish activations to ease INT8 conversion, with the Lite0 variant containing approximately 4.7\,M parameters. MobileViT~\cite{mehta2022mobilevit} introduces hybrid CNN--Transformer blocks; the XXS variant has only 1.3\,M parameters yet is competitive with much larger models on ImageNet. MobileNetV2~\cite{sandler2018mobilenetv2}, included as an off-the-shelf baseline in this work, predates these but remains a strong reference point and exhibits the ReLU6-bounded activation profile we associate with PTQ stability (Section~\ref{sec:results-quant}). We adapt all four as encoders in a U-Net-style decoder.\footnote{The parameter counts quoted here are \emph{backbone-only} (encoder) figures from the classification literature. The full segmentation models we train and deploy add a shared U-Net decoder; their complete encoder$+$decoder parameter counts --- which are the deployment-relevant figures --- are reported in Table~\ref{tab:t1-footprints} (e.g.\ MobileViT-XXS is 1.3\,M as a backbone but 3.09\,M as a full UNet).}

Beyond the four classification-derived backbones we evaluate, several segmentation-specific lightweight architectures have been proposed. BiSeNetV2~\cite{yu2021bisenetv2} uses a two-pathway design (spatial $+$ context) optimised for real-time urban-scene segmentation. PIDNet~\cite{xu2023pidnet} reframes the three-pathway split as a proportional--integral--derivative controller. Fast-SCNN~\cite{poudel2019fastscnn} targets sub-150-MB-FLOP inference at $1024{\times}2048$. We do not evaluate these in this work because (a)~they are designed for and benchmarked on Cityscapes-style urban scenes, not the heterogeneous viewpoints (aerial / oblique / ground) of FSSD, and (b)~our objective is to characterise the deployability of the \emph{commodity} UNet $+$ ImageNet-pretrained-backbone family that the flood-segmentation literature has already converged on, not to propose a new architecture.

\subsection{Post-Training Quantization}
\label{sec:rw-ptq}

Post-training quantization (PTQ) maps FP32 weights and activations to lower-precision integer representations~\cite{jacob2018quantization,krishnamoorthi2018quantizing}. Two main flavours exist: dynamic PTQ, which only quantizes weights and computes activation ranges at inference time, and static PTQ, which calibrates activation ranges offline against a representative dataset. For convolutional networks, the ONNX static QDQ (Quantize--DeQuantize) format with per-channel weight quantization and per-tensor activation quantization is the \emph{de facto} standard~\cite{ortQuant} and is supported across major edge runtimes (ONNX Runtime CPU, TensorFlow Lite, CoreML). The empirical conventional wisdom --- that ``INT8 PTQ retains accuracy within 1 IoU point of FP32''~\cite{krishnamoorthi2018quantizing} --- is derived from classification networks with bounded-range activations (ReLU, ReLU6) on ImageNet-scale benchmarks. As we show in Section~\ref{sec:results-quant}, this assumption does not hold uniformly across the lightweight architectures used in this study: HardSwish-based and attention-based backbones can suffer order-of-magnitude larger PTQ degradation on flood-segmentation pixels than the classification literature would predict.

\subsection{On-Device Disaster Response}
\label{sec:rw-ondevice}

Operational disaster-response platforms must function under intermittent connectivity, limited power, and heterogeneous hardware. Edge AI in humanitarian contexts has been explored by Munawar \emph{et al.}~\cite{munawar2021review} and Kuglitsch \emph{et al.}~\cite{kuglitsch2022facilitating}. The deployment platforms we benchmark --- Apple Silicon CPU and browser WebAssembly --- correspond to two real operational profiles: field-laptop and UAV inference, and citizen-science web portals. Raspberry-Pi-class single-board computers, AWS Graviton ARM CPU, and Android edge targets are discussed in Section~\ref{sec:limitations} as straightforward but unscheduled extensions.
